\pdfoutput=1 
% paper slides: 
% https://docs.google.com/presentation/d/1sa2KAjMVsymqBhlt1iza_Wv2G2EX6TzAd79Tfcalf6I/edit?usp=sharing

% Submitted papers for review should be limited to 12 pages (not counting references). 
%
% Authors of accepted papers will be asked to prepare a final paper of up to 15 pages (including references).
% This is samplepaper.tex, a sample chapter demonstrating the
% LLNCS macro package for Springer Computer Science proceedings;
% Version 2.20 of 2017/10/04
% edit share  https://www.overleaf.com/8724481833ryhbtsyptrtd
\documentclass[runningheads]{llncs}
\usepackage[table,xcdraw,dvipsnames]{xcolor}

\usepackage{graphicx}
\usepackage{tablefootnote}

% Used for displaying a sample figure. If possible, figure files should
% be included in EPS format.
%
% If you use the hyperref package, please uncomment the following line
% to display URLs in blue roman font according to Springer's eBook style:
% \renewcommand\UrlFont{\color{blue}\rmfamily}
\usepackage{todonotes}
%\usepackage[disable]{todonotes}

\usepackage{tabularx}  % Needed for tabularx environment
% \usepackage[font=small]{caption}  % Needed for custom caption
\usepackage[pdfstartview=XYZ,
bookmarks=true,
colorlinks=true,
linkcolor=blue,
urlcolor=blue,
citecolor=blue,
pdftex,
bookmarks=true,
linktocpage=true, % makes the page number as hyperlink in table of content
hyperindex=true
]{hyperref}

\usepackage{multirow}
\usepackage{orcidlink}
\usepackage{textcomp}
% \usepackage{minted}
\usepackage{verbatim}
\usepackage{comment}
\usepackage{threeparttable}
\usepackage{makecell}
\usepackage{xspace}
\usepackage{diagbox}


\newcommand {\ME}[1]{{\textcolor{blue}{(Murali: #1)}}}

\newcommand{\ourwork}{LM4HPC\xspace}
\newcommand{\face}{Hugging Face\xspace}
\newcommand{\crawldata}{OMP4Par\xspace}


%\renewcommand\@makefntext[1]{\parindent 1em\noindent
            %\hspace*{-1em}\@thefnmark\ #1}
            
\begin{document}
%
\title{\ourwork: Towards Effective Language Model Application in High-Performance Computing}
%
\titlerunning{LM4HPC: Towards Effective LM Application in HPC}
% If the paper title is too long for the running head, you can set
% an abbreviated paper title here
%
% \orcidID
\author{Le Chen\inst{1,2}\orcidlink{0000-0003-3847-4108} \and
Pei-Hung Lin\inst{1}\orcidlink{0000-0003-4977-814X} \and
Tristan Vanderbruggen \inst{1}\orcidlink{0000--0000-0000-0000} \and
Chunhua Liao\inst{1}\orcidlink{0000-0001-6477-0547} \and
Murali Emani \inst{3}\orcidlink{0000--0000-0000-0000} \and
Bronis de Supinski\inst{1}\orcidlink{0000-0002-0339-1006}
}
\authorrunning{Chen et al.}
% First names are abbreviated in the running head.
% If there are more than two authors, 'et al.' is used.
%
\institute{Lawrence Livermore National Laboratory, Livermore CA 94550, USA \and
Iowa State University, Ames, IA 50010, USA \and 
Argonne National Laboratory, Lemont, IL 60439, USA}
%
\maketitle              % typeset the header of the contribution


\setlength{\textfloatsep}{8pt}
\begin{abstract}
%The final goal: a success story as follows
In recent years, language models (LMs), such as GPT-4, have been widely used in multiple domains, including natural language processing, visualization, and so on. However, applying them for analyzing and optimizing high-performance computing (HPC) software is still challenging due to the lack of HPC-specific support. In this paper, we design the \ourwork framework to facilitate the research and development of HPC software analyses and optimizations using LMs. Tailored for supporting HPC datasets, AI models, and pipelines, our framework is built on top of a range of components from different levels of the machine learning software stack, with \face-compatible APIs. Using three representative tasks, we evaluated the prototype of our framework. The results show that \ourwork can help users quickly evaluate a set of state-of-the-art models and generate insightful leaderboards.
%The results have been published by a paper accepted by X. 
%\todo[author=Leo,size=\small, color=green]{\textcolor{blue}{CodeBLEU} one or two new metrics should also be included: e.g. edit distance, compilation, execution success rates}
\keywords{Language model  \and Programming language processing \and High-performance computing}
\end{abstract}
%

% 12 pages in total
\section{Introduction} %  (Leo, 1page)
%outline:  1 what are LLMs, why they are important 2. issues to use LLMs for HPC 3. the approach we are using, 4. major findings. 

Language models (LMs) are models designed to understand and generate human language. In recent years, large language models (LLMs) %~\cite{zhao2023survey} 
trained on large amounts of text data have demonstrated stunning capabilities in various natural language processing and visualization tasks. They have also been widely used to process programming languages due to the similarities between natural languages and programming languages.  For example, GPT-4~\cite{bubeck2023sparks} shows early signs of artificial general intelligence. Based on a large language model trained on code~\cite{chen2021evaluating}, GitHub provides an AI assistant for developing software. 

Given the rise of LLMs, it is natural for researchers and developers in the high-performance computing community to start exploiting LMs for addressing various challenges in HPC, including code analysis, code generation, performance optimization, question answering, and so on. However, mainstream frameworks of LMs were originally designed to serve natural language processing. It is difficult for newcomers in HPC to quickly access HPC-specific datasets, models, and pipelines. For example, the current popular \face platform does not include dedicated pipelines for software analyses and optimizations. Another challenge is the entire field is evolving quickly, with new techniques emerging almost weekly, making it challenging for HPC users to keep up with the latest techniques and find relevant ones. Last but not least, there is a lack of standard, reproducible evaluation processes for LMs focusing on HPC-specific tasks. Therefore, it is difficult to have a fair comparison among different models for a given HPC task. 

In this paper, we propose a framework (named \ourwork) designed to serve HPC users as first-class citizens by including internal components and external APIs relevant to HPC-specific tasks. \ourwork's components include models, datasets, pipelines, and so on, while the APIs allow users to interact with these components to finish given HPC tasks. We highlight the contributions of our work as follows:
\begin{itemize}
    \item We design an extensible framework for including and exposing relevant machine learning components to facilitate the adoption of large language models for HPC-specific tasks. 
    \item The framework provides a set of APIs to facilitate essential operations, including code preprocessing, tokenization, integration with new data, and evaluation.
    %\todo{\textcolor{blue}{added in section 4}LangChain is picked to enable size limit challenge. Make sure later text mention this challenge-solution pair}    
    \item A set of pipelines have been developed to support common HPC tasks, including code similarity analysis, parallelism detection, question answering, and so on.   
    \item We provide HPC-specific datasets such as DRB-ML, \crawldata, and OMPQA to support various HPC pipelines.  
    \item Our work introduces standardized workflows and metrics to enable fair and reproducible evaluation of LLMs for HPC-specific tasks.
    \item Using three representative tasks, we demonstrated how the framework can be used to test a set of language models and generate leaderboards.
\end{itemize}

%The remainder of this paper is organized as follows. In the next section, we introduce the background of large language models and their applications in multiple domains. Section~\ref{sec-approach} presents the design of the LM4HPC framework, including its internal components and external APIs. Section~\ref{sec-evaluation} gives a preliminary results of the current implementation, using three representative HPC tasks to generate leaderboards of some popular models. Related work is presented in Section~\ref{sec-related-work}. Finally, Section~\ref{sec-conclusion} concludes the work and discusses future work. 

% Leo

% \section{Background: Pei-Hung - 2}
% \label{sec-background}
% \input{sections/2-background}

\section{Background} % 2 pages?
\label{sec-background}
%
% 1.1 what is a language model
% 1.2 introduction of transformers
% 1.3 large language models
% 2.1 transformers based models for code analysis, such as codebert, codet5
% 2.2 large language model for code analysis and HPC analysis
% 1.4 what is missing for LM or LLM for HPC
%\todo{A figure about large language model may be needed to help explain its components}


Language models (LMs) are machine learning models designed to comprehend and generate human language. They can be used to facilitate natural and intuitive interactions between humans and machines. 
Early generations of LMs, using recurrent neural networks (RNNs),  showed inspiring results for various natural language processing (NLP) tasks.  
A transformative evolution by the Transformer \cite{vaswani2017attention} reveals remarkable potentials of LMs.
Introduced by Vaswani \textit{et al.}, transformer
models utilize the attention mechanism to capture the dependencies between all words in an input sentence, irrespective of their positions. Compared to RNNs, transformers process data in parallel rather than sequentially and significantly improve the efficiency of model training and inference. %Transformer architecture has become a cornerstone in the field of NLP. 
Transformers further enables the inauguration of the large language models (LLMs).
%Large language models (LLMs) are a type of language models that are trained on a vast amount of data. 
Compared to LMs, LLMs are trained on a vast amount of data and possess parameter counts on the order of billions or more, allowing them to generate more detailed and nuanced responses. Examples of LLMs include OpenAI's GPT-3, GPT-4 and Google's BARD. Nowadays, LLMs have shown remarkable capabilities in NLP tasks like translation, question answering, and text generation.

%The evolution of language models has been transformative, migrating from recurrent neural networks (RNNs) to the more advanced transformer models~\cite{vaswani2017attention}, with remarkable success across various natural language processing (NLP) tasks. 
% However, this venture comes with its unique set of challenges:


% \begin{table}[htbp]
% \centering
% \caption{Language models, associated training data and tasks}
% \label{tab-modelExample}
% \begin{threeparttable}
% \begin{tabular}{|c|c|c|c|c|c|}
% \hline
% \multicolumn{2}{|c}{\textbf{Model}} &
% \multicolumn{2}{|c|}{\textbf{Training Data}} &
% \multirow[t]{2}{*}{\textbf{TokenLimit}} &
% \multirow[t]{2}{*}{\textbf{Availability}} \\
% \cline{1-4}
%  \textbf{Name} & \textbf{Size} & \textbf{Type} & \textbf{Size} &  & \\ \hline
%  BERT & 340M & Text & 3.5B words & 512 & code, weights\\ \hline
%  CodeBERT & 125M & Mixed & \makecell{2.1M (bimodal)\\6.4M (unimodal)} & 512 & code, weights\\ \hline
%  GraphCodeBERT & 110M & Code & 2.3M functions  & 512 & code, weights\\ \hline
%  CodeT5 & 770M & Code &  8.35M instances  & 512 & code, weights\\ \hline
%  GPT-3 & 175B & Mixed &  500B words  & 4096 & API\\ \hline
% Cerebras-GPT & 0.11B$\sim$13B & Mixed &  800GB  & 2048 & code, weights\\ \hline
% LLaMA & 7$\sim$65B & Mixed &  1.4T tokens  & 4096 & code\\ \hline
% StarCoder & 15B & Code &  1T tokens  & 8192  & code, weights\\ \hline
% Dolly 2.0 & 3$\sim$12B & Text & 15k inst./response  & 2048 & code, weights\\ \hline
% BARD & 1.6B & Mixed & 1.56T words & 1000 & API\\ \hline
% GPT-4 & $>$ 1T & Mixed & undisclosed & 8k/32k & API\\ \hline
% \end{tabular}
% %token limits for gpt-4 and gpt-4-32k are 8192 and 32768 respectively.


% %\begin{tablenotes}\footnotesize
% %\item[*]NLP: natural language processing, PLP: programming language processing \newline
% %LM: language model, LLM: large language model \newline
% %NLT: natural language text, PLT: programming language text \newline

% %\end{tablenotes}
% \end{threeparttable}

% \end{table}



% Please add the following required packages to your document preamble:
% \usepackage{multirow}
% Please add the following required packages to your document preamble:
% \usepackage{multirow}
% Please add the following required packages to your document preamble:
% \usepackage{multirow}

%table in a google sheet: 
%%%  https://docs.google.com/spreadsheets/d/1oS3aB7QsL7dkcCm6Lkv_IviBrFzlKiw_3VM0V8LaNOM/edit?usp=sharing
% Please add the following required packages to your document preamble:
% \usepackage{multirow}
%\todo{\textcolor{blue}{added in table 1} add  NVIDIA’s Megatron: }
\vspace{-12pt}
\begin{table}[]
\caption{Language models, associated training data and tasks}
\label{tab-modelExample}
\begin{tabular}{|ccc|cc|c|c|}
\hline
\multicolumn{3}{|c|}{\textbf{Model}} &
  \multicolumn{2}{c|}{\textbf{Training data}} &
  \multirow{2}{*}{\textbf{\begin{tabular}[c]{@{}c@{}}Token\\  Limit\end{tabular}}} &
  \multirow{2}{*}{\textbf{Avail.}} \\ \cline{1-5}
\multicolumn{1}{|c|}{\textbf{Name}} &
  \multicolumn{1}{c|}{\textbf{Released}} &
  \textbf{Size} &
  \multicolumn{1}{c|}{\textbf{Type}} &
  \textbf{Size} &
   &
   \\ \hline
\multicolumn{1}{|c|}{BERT}                     & \multicolumn{1}{c|}{2018/10}       & 340M      & \multicolumn{1}{c|}{Text}  & 3.5B words      & 512    & Weights \\ \hline
\multicolumn{1}{|c|}{CodeBERT} &
  \multicolumn{1}{c|}{2020/11} &
  125M &
  \multicolumn{1}{c|}{Mixed} &
  \begin{tabular}[c]{@{}c@{}}2.1M(bimodal)\\  6.4M (unimodal)\end{tabular} &
  512 &
  Weights \\ \hline
  \multicolumn{1}{|c|}{Megatron}            & \multicolumn{1}{c|}{2021/04}       & 1T      & \multicolumn{1}{c|}{Text}  & 174 GB  & 512    & Weights \\ \hline
\multicolumn{1}{|c|}{GraphCodeBERT}            & \multicolumn{1}{c|}{2021/05}       & 110M      & \multicolumn{1}{c|}{Code}  & 2.3M functions  & 512    & Weights \\ \hline
\multicolumn{1}{|c|}{CodeT5}                   & \multicolumn{1}{c|}{2021/11}       & 770M      & \multicolumn{1}{c|}{Code}  & 8.35M instances & 512    & Weights \\ \hline
\multicolumn{1}{|c|}{GPT-3}                    & \multicolumn{1}{c|}{2022/03/15} & 175B      & \multicolumn{1}{c|}{Mixed} & 500B tokens           & 4096   & Weights \\ \hline
\multicolumn{1}{|c|}{LLaMA}                    & \multicolumn{1}{c|}{2023/02/24}   & 7$\sim$65B     & \multicolumn{1}{c|}{Mixed} & 1.4T tokens      & 4096   & Weights*          \\ \hline
\multicolumn{1}{|c|}{GPT-4}                    & \multicolumn{1}{c|}{2023/03/14}   & 1T        & \multicolumn{1}{c|}{Mixed} & undisclosed     & 8k/32k & API*           \\ \hline
\multicolumn{1}{|c|}{BARD}                     & \multicolumn{1}{c|}{2023/03/21}   & 1.6B      & \multicolumn{1}{c|}{Mixed} & 1.56T words            & 1000   & API           \\ \hline
\multicolumn{1}{|c|}{Cerebras-GPT}             & \multicolumn{1}{c|}{2023/03/28}    & 0.11$\sim$13B & \multicolumn{1}{c|}{Text} & 800 GB           & 2048   & Weights \\ \hline
\multicolumn{1}{|c|}{Dolly 2.0}                & \multicolumn{1}{c|}{2023/04/12}   & 3$\sim$12B     & \multicolumn{1}{c|}{Text}  & 15k instr./resp.             & 2048   & Weights \\ \hline
\multicolumn{1}{|c|}{\begin{tabular}[c]{@{}c@{}}StarCoder\end{tabular}} & \multicolumn{1}{c|}{2023/05/4}   & 15B       & \multicolumn{1}{c|}{Code}  & 1T tokens            & 8192   & Weights \\ \hline
\multicolumn{1}{|c|}{\begin{tabular}[c]{@{}c@{}} StarChat-Alpha\end{tabular}} & \multicolumn{1}{c|}{2023/05/4}   & 16B       & \multicolumn{1}{c|}{Code}  & 31k instr./resp.            & 8192   & Weights \\ \hline
\end{tabular}
\end{table}
\vspace{-15pt}

%\subsection{Transformers and LLMs}
%\subsection{LM/LLM Tasks}
% \vspace{-10pt}
Table~\ref{tab-modelExample} shows some example language models and their release dates, sizes, training data, input token length limits, and availability. LLaMA~\cite{touvron2023llama}'s weights can be obtained after filling out some form. GPT-4 has a waiting list to use its API. At the time of writing this paper, we have not yet obtained its access.
 
%Deploying LM applications could be complex as it requires an understanding of datasets, machine learning frameworks, models, and other technical aspects. This complexity is further amplified by the rapid evolution and expansion of LLMs. Hugging Face has become the leading platform for the machine-learning community to share and explore cutting-edge models and datasets. It supports all the mainstream machine-learning frameworks and offers thoughtfully designed APIs for the effective development of LM applications.  Hugging Face simplifies the process of deploying sophisticated LMs, making state-of-the-art NLP technology more accessible and manageable for developers and researchers.

% \textit{at least one billion or more parameters.}}
%\subsection{LM and LLM for Code}
LMs are trained mainly by text data with a focus on NLP.  The sources of the training data mainly come from books, web content, newspapers, scientific articles, and  other text data in various natural languages.    
 Latest LLMs have demonstrated rich skill sets in NLP including text prediction, common sense reasoning, reading comprehension, translation and question answering.

 
There has been a keen interest in deploying NLP techniques to programming language processing (PLP) tasks, such as code summarization, code generation, and code similarity analysis \cite{chen2022multi,flynn2022finding}. 
Previous studies have demonstrated successful applications of traditional language models to PLP tasks, showing the feasibility of this approach~\cite{devlin2018bert}.
%Previous research has extended transformer-based language models to the PLP domain. 
CodeBERT \cite{feng2020codebert}, for example, is a transformer-based model trained with a diverse range of programming languages and can be used for a variety of programming-related tasks. Similarly, CodeT5 \cite{wang2021codet5} is a variant of Google's T5 language model, trained specifically on code datasets to perform advanced programming tasks like code completion, bug detection, and code summarization.
Lately, StarCoder \cite{li2023starcoder}, a 15B parameter model trained with 1 trillion tokens sourced from a large collection of permissively licensed GitHub repositories, is developed to be a Large Language Model mainly for code generation or completion. 
StarChat-Alpha is a GPT-like chat model fine-tuned from StarCoder to act as a helpful coding assistant. 

% \color{red}{add starcoder}
% \textcolor{blue}{add starcoder; add HPCFAIR}
\vspace{-10pt}
\subsection{LMs for HPC}
With the recent breakthroughs in Generative Pretrained Transformer (GPT) large language models \cite{brown2020language}, it has become increasingly intriguing to explore the application of large language models (LLMs) for HPC tasks.
%Despite the significant advances made by Large Language Models (LLMs) in comprehending and generating both natural and programming languages, 
However, their deployment in the HPC domain is still relatively unexplored. This venture comes with various challenges, including:
\vspace{-4pt}

%\todo{\textcolor{blue}{added example} although there are a lot HPC code in ..., but not reprocessed ...}
\begin{enumerate}
    \item Pipelines: Traditional language model frameworks like Hugging Face were designed to support natural language processing or compute vision problems. Expanding LMs to any new domain, including HPC,  requires the addition of new pipelines designed to finish domain-specific tasks.   
    \item Datasets: The HPC domain encompasses an extensive amount of code spanning various fields, including biology and climate modeling. However, preparing this data for machine learning training, such as labeling parallelizable loops in HPC programs for parallelism detection, presents significant challenges. The scarcity of ready-to-use, pre-labeled HPC datasets poses a particular obstacle for training language models, especially large ones, highlighting the need for more shared resources in the community.
    \item Pre-processing: Pre-processing in the context of LMs for HPC typically involves the conversion of source files into a sequence of tokens. However, the direct application of NLP tokenizers to code can be sub-optimal. For instance, an NLP tokenizer might split a variable name into two tokens, a scenario that is not desirable for PLP analysis.  Also, models designed for processing source code may use graph representations, such as abstract syntax trees, to have better performance.
    \item Input size limit: Language models often have limited input token lengths (such as 512 to a few thousand of tokens). HPC tasks often involve processing large-scale software packages with millions of lines of source code. 
%    \item Fine-tuning:  To fully harness the capabilities of state-of-the-art LLMs, fine-tuning these models with HPC-specific data is necessary. This process demands a thorough understanding and expertise in handling LLMs. 
    \item Evaluation: There is a pressing need for standardized and reproducible evaluation of different models in the context of various HPC tasks, using metrics suitable for domain-specific requirements. 
    %This is particularly challenging given the stochastic nature of language models. The resulting leaderboards can facilitate communication regarding both the state-of-the-art and areas that are lacking. 
\end{enumerate}
% \vspace{-10pt}



% Pei-Hung
\vspace{-14pt}
\section{Approach} %3.5 pages
\label{sec-approach}
% \vspace{-10pt}
To address the challenges discussed in Section~\ref{sec-background}, we introduce \ourwork, a comprehensive framework that encapsulates a suite of machine learning components within user-friendly APIs. This framework is tailored for HPC users, simplifying the implementation process and making the robust capabilities of language models more accessible and user-friendly within the HPC community. The primary goal of \ourwork is to reduce the complexities inherent in employing language models, thus enabling HPC users to leverage their powerful capabilities more effectively and efficiently.




\subsection{\ourwork  design overview}
% DO NOT ADVERTIZE HUGGING FACE IN our own design overview
%While Hugging Face offers maturing designs for each step in a machine-learning application, a gap remains for HPC users looking to leverage these benefits. 
%To bridge this gap, we introduce \ourwork, a scalable and extensible framework specifically designed to effectively deploy language models for the HPC domain. 
%In line with the FAIR principles, \ourwork is built atop Hugging Face. This approach enables HPC users to easily access and deploy models and datasets shared by the NLP community. The APIs of Hugging Face and their documentation help HPC users reduce the steep learning curve typically associated with language models.

Figure~\ref{fig:architecture} provides the overview of the \ourwork framework. It is built on top of multiple internal machine learning components with \face-compatible APIs. Higher-level components provide concepts and interfaces to users, while middle or lower-level components provide implementation support. 
Table \ref{tab:api_table} shows the example classes and functions in LM4HPC API, including those supporting HPC-specific language models,  tokenizers for programming languages, datasets, inference pipelines, and evaluation.
%Figure~\ref{fig:architecture} provides an overview of the \ourwork framework. Delivered as a Python package, \ourwork aims to simplify the deployment of large language models within the HPC domain. The framework employs \face as the underlying runtime library to communicate between machine learning frameworks and higher-level APIs.
%The higher-level components within \ourwork comprise five primary parts: dataset, tokenizer, pipelines, vector stores, and evaluations. 
We elaborate on some essential components in the following subsections. 

\begin{figure}[h]
    \centering
    \includegraphics[width=\textwidth]{images/architecture.pdf}
    \caption{Overview of the \ourwork framework}
    \label{fig:architecture}
\end{figure}



% Please add the following required packages to your document preamble:
% \usepackage{multirow}
% \usepackage[table,xcdraw]{xcolor}
% If you use beamer only pass "xcolor=table" option, i.e. \documentclass[xcolor=table]{beamer}
\begin{table}[!t]
\centering
\caption{LM4HPC API: example classes and functions. Each class can be imported using ``from lm4hpc import *'' in Python.}
\label{tab:api_table}
\footnotesize 
\setlength{\tabcolsep}{2pt} 
\begin{tabular}{|p{0.13\columnwidth}|p{0.25\columnwidth}|p{0.56\columnwidth}|}
\hline
\textbf{\scriptsize LM4HPC classes} & \textbf{\scriptsize Description} & \textbf{\scriptsize Example API functions} \\
\hline
\multirow{3}{\linewidth}{\scriptsize \texttt{hpcmodel}} &
\multirow{3}{\linewidth}{\raggedright\arraybackslash \scriptsize Fine-tune text-based (HF, OpenAI) and graph-based models, including local private ones, for HPC tasks} &
\raggedright\arraybackslash \scriptsize \texttt{hpcmodel.from\_pretrained(model\_name\_or\_path: Optional[str], *model\_args, **kwargs)} \\ \cline{3-3} 
& & \raggedright\arraybackslash \scriptsize \texttt{hpcmodel.save\_pretrained(model\_name\_or\_path: str, *model\_args, **kwargs)} \\ \cline{3-3} 
& & \raggedright\arraybackslash \scriptsize \texttt{hpcmodel.finetune()} \\
\hline
\multirow{3}{\linewidth}{\scriptsize \texttt{hpctokenizer}} &
\multirow{3}{\linewidth}{\raggedright\arraybackslash \scriptsize APIs to represent code in either tokenized text, trees, or graphs} &
\raggedright\arraybackslash \scriptsize \texttt{hpctokenizer.from\_pretrained(model\_name\_or\_path: Optional[str], *model\_args, **kwargs)} \\ \cline{3-3} 
& & \raggedright\arraybackslash \scriptsize \texttt{hpctokenizer.addtokens(contentsingle\_word=False, strip=False, normalized=True)} \\ \cline{3-3} 
& & \raggedright\arraybackslash \scriptsize \texttt{hpctokenizer.encoding()} \\
\hline
\multirow{4}{\linewidth}{\scriptsize \texttt{hpcdatasets}} &
\multirow{4}{\linewidth}{\raggedright\arraybackslash \scriptsize Load and process HPC datasets} &
\raggedright\arraybackslash \scriptsize \texttt{hpcdatasets.load(path: str, data\_files: Union[str, List, Dict, None], **kwargs)} \\ \cline{3-3} 
& & \raggedright\arraybackslash \scriptsize \texttt{hpcdatasets.split(dataset: hpcdatasets, partition: [float, float, float], **kwargs)} \\ \cline{3-3} 
& & \raggedright\arraybackslash \scriptsize \texttt{hpcdatasets.shuffle(dataset: hpcdatasets, **kwargs)} \\ \cline{3-3} 
& & \raggedright\arraybackslash \scriptsize \texttt{hpcdatasets.sort(dataset: hpcdatasets, **kwargs)} \\
\hline
\scriptsize \texttt{hpcpipeline} &
\scriptsize Pre-built pipelines for common PLP tasks &
\raggedright\arraybackslash \scriptsize \texttt{hpcpipeline(task: str, model\_name\_or\_path: str, *model\_args, **kwargs)} \\
\hline
\multirow{2}{\linewidth}{\scriptsize \texttt{hpceval}} &
\multirow{2}{\linewidth}{\raggedright\arraybackslash \scriptsize Evaluate the results of various models} &
\raggedright\arraybackslash \scriptsize \texttt{hpceval.compute(task: str, models\_name\_or\_path: [[str]], data\_files: Union[str, List, Dict, None], *model\_args, **kwargs)} \\ \cline{3-3} 
& & \raggedright\arraybackslash \scriptsize \texttt{hpceval.plot(shape: str)} \\
\hline
\end{tabular}
\end{table}


% \vspace{-15pt}
%========================================
\subsection{HPC Tasks and Inference Pipelines}
HPC users are interested in a wide range of tasks related to programming language processing. Table~\ref{tab-HPC-tasks} outlines one way to categorize HPC-specific tasks. The purpose here is not to provide a comprehensive taxonomy of all tasks but a starting point for common tasks we are interested in supporting in our framework. Most tasks are self-explanatory by names and each may have further sub-tasks. For example, clone detection can be viewed as a specialized sub-task under code similarity analysis. 
\begin{table}[htbp]
\centering
\caption{HPC Tasks for Programming Language Processing: Categories and Examples}
\label{tab-HPC-tasks}
\begin{tabular}{l|l|l}
\hline
%\multicolumn{3}{c}{\textbf{Category}} \\ \hline
\textbf{Code Analysis} & \textbf{Code Generation} & \textbf{Others} \\ \hline
Compiler Analysis & Code Completion & Test Case Generation \\
Algorithm Classification & Natural Language-to-Code & Code Search \\
Code Similarity Analysis & Code Translation & Question Answering \\
Documentation Generation & Code Repair & Code Review \\
Parallelism Detection & Code Migration & Decompilation  \\
Defect Detection & Compilation &  IR-to-Source Translation \\ \hline
\end{tabular}
\end{table}

In the context of machine learning, a pipeline represents a sequence of data processing stages to complete a task. Our \ourwork framework extends the pipeline function provided by \face, adapting it for HPC tasks.  
%Pipelines allow HPC researchers and developers, even those with limited ML experience, to rapidly harness the power of cutting-edge LLMs. Currently, \face provides pipelines for NLP and image processing tasks only. 
%\todo{HPCFAIR vs. HPC-FAIR: which one do you mean?  This paper focuses on HPC-FAIR. We only need to briefly mention previous HPCFAIR in related work. }
%In an effort to streamline this process and provide user-friendly interfaces for executing a variety of HPC-related inference tasks using pre-trained models, we have extended the \face pipeline() API within \ourwork. 
We have developed three inference pipelines: code similarity analysis, parallelism detection, and OpenMP question answering. Code similarity analysis determines the similarity between a pair of code snippets. 
%The exact meaning of similarity depends on the context. It could mean similar syntax, semantics, or functionality. 
Parallel detection is defined to check if an input code snippet can be parallelized or not using OpenMP.  The OpenMP question answering pipeline is designed to use models to generate answers to OpenMP-related questions. %\todo{Double check if the task description is accurate} %These pipelines are designed to significantly reduce the complexity of integrating machine learning models into HPC workflows.
% \begin{comment}  % No space   

% Listing \ref{ls:pipepline} illustrates how to carry out inference tasks using \ourwork. In line 3, the parallelism detection pipeline is loaded by passing the pipeline identifier to the \mintinline{python}{pipeline()} API. 
% With just two lines of code, the pipeline API can generate results for parallelism detection tasks. Each pipeline task is equipped with predefined pre-trained models and tokenizers. This efficient design drastically reduces the complexity of these tasks, enabling users to obtain and evaluate results quickly and effortlessly.

% \begin{listing}[h]
% \caption{Example of using \ourwork pipeline.}
% \label{ls:pipepline}
% \begin{minted}
% [
% frame=lines,
% framesep=2mm,
% baselinestretch=1.2,
% % bgcolor=LightGray,
% fontsize=\footnotesize,
% linenos
% ]{python}
% from \ourwork import pipeline
% # parallelism detection task
% pipe = pipeline("parallelism-detection")
% res = pipeline(code)
% \end{minted}
% \end{listing}
% \end{comment}
% we have to mention pipelines before Datasets. Otherwise hard to explain datasets internals like enabling parallel detection, similarity analysis etc.

Tokenizers are responsible for preprocessing input into an array of numbers as inputs to a model. They are essential components used by pipelines. Most LM tokenizers are primarily designed for NLP tasks. 
%Consequently, they often fall short when directly applied to HPC tasks. 
For instance, given a function name \texttt{my\_func}, a typical NLP tokenizer like BERT might split it into separate tokens (such as `\texttt{my}', `\texttt{\_}', and `\texttt{func}') while a code-aware tokenizer may treat the function name as a single entity to ensure a more meaningful representation. % of the code. %This clearly demonstrates the limitations of NLP tokenizers in correctly parsing code.

To overcome this, we developed the \ourwork tokenizer, leveraging the treesitter~\cite{treesitter} and programl~\cite{programl} library. Our tokenizer is specifically designed to handle the pre-processing of code data required for a language model. It includes tokenizers such as the ast-tokenizer.
%which parses code into tokens based on Abstract Syntax Tree (AST) traversal, thereby ensuring a more accurate representation of the code. 
As a result, \ourwork can accommodate models (such as augAST~\cite{chen2023learning-dataset}) that require AST as input in the pipeline. % and DeepSCC~\cite{}.
%augAST 

\subsection{Datasets}
\label{sec-datasets}
Datasets are crucial for any machine learning application. 
%The NLP community shares datasets via \face, establishing a platform to evaluate the performance of different models using the same dataset. 
Within the \ourwork framework, we contribute HPC-specific datasets either by converting existing ones into \face-compatible formats or by creating new ones from scratch. 

%in two significant ways: %1. We convert existing HPC datasets into a \face compatible format and share them on the \face Hub. This is achieved by transforming the HPC data into a JSON file where each row represents a single example from the HPC dataset. 
% 2. We create new HPC datasets specifically designed for LLMs. In this work, we have processed three datasets and evaluated different models with these datasets in section 4. 
    %1) We convert existing HPC datasets into a \face-compatible format and share them on the \face Hub. This is achieved by transforming the HPC data into a JSON file where each row represents a single example from the HPC dataset. 
    %2) We create new HPC datasets.     
%In this work, we have processed three datasets and evaluated different models with these datasets in section 4. 


We have converted three existing datasets to be compatible with \face dataset API: POJ-104~\cite{dataset:poj104}, DRB-ML~\cite{osti_1958879}, and \crawldata~\cite{chen2023learning-dataset}.  
POJ-104 is derived from a pedagogical programming open judge (OJ) system that automatically evaluates the validity of submitted source code for specific problems by executing the code. This dataset is particularly useful for the code similarity task. % as it consists of 104 problems, each accompanied by 500 student-written C/C++ programs that provide solutions to the respective problem. 
The DRB-ML dataset contains 658 C/C++ OpenMP kernels derived from DataRaceBench~\cite{dataset:dataracebench}. We extended it to have labels indicating if a kernel is parallelizable or not. 
The \crawldata dataset is an open-source benchmark composed of data from three resources: code crawled from GitHub, OpenMP benchmarks such as Nas Parallel Benchmarks~\cite{jin1999openmp} and Rodinia~\cite{che2009rodinia}, and synthetic code. This dataset contains loops with labels indicating whether a loop is parallel and, if parallelizable, the corresponding OpenMP directive associated with the loop. 

Furthermore, we have manually created a new OpenMP question answering dataset called OMPQA in order to probe the capabilities of language models in single-turn interactions with users. Similar to other QA datasets, we include some request-response pairs which are not strictly question-answering pairs. 
%We manually 
%aim to create a comprehensive, diverse, and balanced OpenMP Q\&A dataset. 
The categories and examples of questions in the OMPQA dataset can be found in  Table~\ref{tab-omp-qa}. 
%The basics category include question-answer pairs related to basic OpenMP concepts. The example category collects questions about code examples written in OpenMP. The compilers category list questions regarding OpenMP compilers, including their usage and internals. The benchmarks category has questions probing the knowledge about popular OpenMP benchmarks for performance, correctness, and so on.


\begin{table}[h]
\centering
\begin{tabular}{|c|c|p{3.7in}|}
\hline
 \textbf{Category}&  \textbf{Count} & \textbf{Example Questions}    \\
\hline
 Basics & 40 &  What is a worksharing construct in OpenMP? \\
\hline
Examples & 20 & Give an example OpenMP C code for computing PI using numerical integration. \\
\hline
Compilers & 24 &  {In what language is LLVM written? \newline How is a parallel region represented in Clang?} \\
\hline
 Benchmarks & 23 & {What are the NAS Parallel benchmarks? \newline Which benchmark assesses data race detection tools?}  \\
\hline 
\end{tabular}
\caption{OMPQA: categories and examples of questions}
% current collection of Q&A
% https://docs.google.com/spreadsheets/d/1dqsd2v_lcRyLv8XiLIy00E_ezBiaivcSjZrMDqFeu_w/edit#gid=0
\label{tab-omp-qa}
\end{table}


\begin{comment}
    
%\subsection{Tokenizer} % merged into the discussion of pipelines
Most LM tokenizers are primarily designed for NLP tasks. Consequently, they often fall short when directly applied to HPC tasks. For instance, given a function name \mintinline[]{python}{my_func}, an NLP tokenizer such as BERT might split it into separate tokens [\mintinline[breaklines, breakafter=]{bash}{'my', '_', 'func'}]. This clearly demonstrates the limitations of NLP tokenizers in correctly parsing code.

To address this issue, we have developed the \ourwork tokenizer, built on top of various code tokenizer libraries. The \ourwork tokenizer is specifically designed to handle all the pre-processing of code data required for a language model. It includes tokenizers such as the ast-tokenizer, which parses code into tokens based on Abstract Syntax Tree (AST) traversal, thereby ensuring a more accurate representation of the code.

Listing \ref{ls:hpcfair-tokenizer} \todo{Where is the listing? We don't have much space, we may just briefly mention tokenizers  as part of the pipeline}demonstrates the steps involved in tokenizing code and processing it through models. This procedure can be complex for those with limited machine learning backgrounds. Challenges such as selecting the appropriate model, choosing the right tokenizer, and managing post-processing may demand substantial effort to address. 
\end{comment}

\vspace{-25pt}
\subsection{Integration With New Data}
% https://drive.google.com/file/d/17iV8BbkY5QCwT3G6r6zfwQaOpYanWjV4/view?usp=sharing
% LangChain source code

% https://docs.google.com/drawings/d/10c3a9qfFxfLctr9KiROoK2xy5NhLFfW2TMg3cr_S-x4/edit?usp=sharing
% langchain figure if space allows
%Large language models are trained on a variety of data and have demonstrated their competence in a wide range of NLP tasks. Remarkably, they can also handle some HPC tasks when their training data contains programming languages. 
%However, these out-of-the-box models may not always provide the optimal performance for more complex HPC tasks. 
Language models derive knowledge from training datasets and store this knowledge in internal weights within the model's neural network architecture. However, incorporating new information into a trained model presents a challenge. Traditionally, one might fine-tune pre-trained models with their own data for specific tasks, but this approach requires substantial relevant data and can be resource-intensive. An alternative approach involves integrating new data as context information into a user prompt, but this is constrained by the limited input token lengths of current models.

To address this challenge, \ourwork leverages the LangChain framework~\cite{langchain2023} to easily integrate new data. LangChain aggregates a wide variety of components to build applications using LLMs. Particularly, it provides APIs allowing LLM applications to store large amounts of text in semantic databases called vector stores. 
The way to integrate new data can be done in two steps. First, text data is chunked and embedded with an LLM before being saved into a vector store. 
%Later, The embeddings and the corresponding text chunks are then retrieved using similarity in the latent space. 
Later, user prompts are matched with relevant chunks in the vector store using similarity analysis. 
%The matching chunks are injected in the prompt to provide context.
The top-matched chunks are then injected into the original prompts to form a new prompt with relevant context information. By employing this new prompt, language models can generate answers that incorporate new and relevant user data while still staying within the token length limits.
%generate Vector Store from their codebase or documentation. Initial queries from users are sent to the Vector Store for a similarity search instead of being sent directly to the LLMs. Consequently, LLMs receive both the initial query and the relevant information from the users' database, enabling them to generate more accurate answers or actions.
%Our solution with Langchain vector storage is designed to integrate user data with LLMs, offering a more convenient and efficient approach to tailor the LLMs to the specifics of users' HPC tasks, thereby enhancing their performance.

\vspace{-10pt}
\subsection{Evaluation}
An easily accessible harness for evaluating different language models on HPC tasks is crucial. Standard and reproducible results from such evaluations can provide researchers and developers with insightful starting points, helping them select suitable models for their specific needs and identify research opportunities. 

In response to this need, we developed an evaluator API in \ourwork. One challenge we encountered is the lack of standardized metrics for code evaluation. Unlike natural language tasks, where metrics such as BLEU,
% ~\cite{papineni2002bleu}, 
ROUGE,
% ~\cite{lin2004rouge}
and METEOR
% ~\cite{banerjee2005meteor} 
are commonly used, the domain of code lacks such universally accepted measures of quality. We are adding various LLM metrics such as CodeBLEU~\cite{ren2020codebleu} for code output.
%These metrics aim to provide a more comprehensive and accurate evaluation of models' performance on code-related tasks. 
Another challenge is that language models may generate different answers for the same input in different inference runs. Evaluation should consider consistent sampling settings (such as temperatures) and control over random seeds to improve reproducibility.  


%By integrating these metrics into \ourwork, we hope to facilitate more robust, effective, and straightforward model evaluation for HPC tasks, ultimately aiding in the selection and development of more advanced models for these complex tasks.
%\todo{\textcolor{blue}{CodeBLEU} what new metrics?}

%\todo{\textcolor{blue}{added }mention that we create and release a set of test harnesses (scripts) to enable standard and reproducible evaluation of supported HPC tasks.}

Ultimately, many users are interested in seeing leaderboards that showcase mainstream models competing on common HPC tasks. To satisfy this interest, we create and release a set of test harnesses scripts to enable standard and reproducible evaluation for supported HPC tasks.
% With its comprehensive components and APIs, our framework supports the generation of such leaderboards using a set of test harness scripts. 
%This allows the community to easily identify the most appropriate models for a given task, as well as the areas that require further research and development.  

%\todo{\textcolor{blue}{table 2} We miss the table of the APIs we provide and what functionality each of them provides}


\vspace{-10pt}

% \section{Approach: Le - 3.5}
% \input{sections/3-approach}
% \label{sec-approach}
% Le

%\todo{Please consider rename "Evaluation" section to be "Preliminary Results". Then we don't have name collision to the evaluation of models in our framework. In the framework, the term "evaluation" is already used to mean the test harness benchmarking one or more models for a given task, using some datasets, optionally generating leaderboards.  \textcolor{blue}{We have changed the section name to "preliminary results". The word evaluation is used in the approach section}}

\section{Preliminary Results} %  Le - 3.5
\label{sec-evaluation}
In this section, we evaluate the current prototype implementation of LM4HPC through experiments designed to generate leaderboards for three representative tasks: Code Similarity Analysis, Parallelism Detection, and OpenMP Question Answering. \ourwork utilizes LangChain v0.0.174, \face's transformers v4.29.0 and datasets v2.12.0 as our runtime libraries. 
%The models and datasets used for evaluation are loaded from the local SSD drive or \face . 
Details of the models and datasets will be discussed in subsequent subsections.

%\subsection{Hardware and Software Configuration}
%In this section, we conduct evaluations of the results from HPC inference tasks, generated using HPC-FAIR APIs. 
Our experiments were conducted on two machines: 1) a Google Colab VM with a 6-core Xeon processor operating at 2.20GHz, 83.5 GB main memory, 166GB HDD drive, and an NVIDIA A100 GPU with 40 GB memory. 2) a Dell workstation equipped with a dual Intel Xeon 6238 CPU operating at 2.10GHz, 128 GB main memory, 1TB SSD drive, and an NVIDIA Quadro RTX 6000 GPU with 24GB memory.  The majority of our experiments were run on the Google Colab machine to leverage its superior GPU memory. However, we encountered difficulties running Cerebras-GPT on the Colab machine and were compelled to use the Dell workstation with larger CPU memory instead. 
 % The system runs on Red Hat Enterprise Linux 8 OS with CUDA 12.0 installed.  
 
%\todo{\textcolor{blue}{added google colab }faithfully list all machines we used and why a single machine cannot run them all. This can be mentioned in conclusion as a challenge to run models}
%\todo{\textcolor{blue}{added StarCoder; we don't have time for LLaMA }add two more models: StarCoder and LLaMA into the results tables}



% For specialized tasks, such as parallelism detection, out-of-the-box models typically underperform. To address this, we have integrated models from previous research into HPC-FAIR. We will discuss these additions in the following sections.
%\todo{Mention model parameter sizes, temperature used, memory footprints, etc. in the table or text}

% \todo[author=Le, color=green]{add model parameters for each model; input and output of the model}
% \todo{Merge some large language models of this table into Table 1. This table should only contain additional smaller models for evaluation.}
% \begin{table}[h]
% \caption{Models used in the evaluation. For LLM model, we mark the parameter size used for evaluation. }
% \label{tab:eval-models}
% \begin{tabular}{c|cccc}
% \hline
% Model Type           & Model Name   & temperature & application             & pre-trained with code data? \\ \hline
% \multirow{2}{*}{LM}  & Codebert    &  & code specified & Yes                         \\ \cline{2-4} 
%                      & GraphCodeBert    & & code specified    & Yes                         \\ \hline
% \multirow{3}{*}{LLM} & ChatGPT 3.5 - 175B &0.7 & general text generation & Yes                         \\ \cline{2-4} 
%                      & LLaMA - 13B   &     & general text generation & No                          \\ \cline{2-4} 
%                      & Dolly 2.0 - 12B  &  & general text generation & No                          \\ \hline
% \end{tabular}
% \end{table}

% \subsection{Datasets}
% \todo{In the approach section, we also have the datasets subsection to introduce HPC specific datasets we are adding. So most text about our own datasets should be moved there. This Preliminary result section should be very brief, just introducing other datasets and referring to our own datasets. }
% In this section, we introduce the datasets used in the evaluation experiments for our HPC-FAIR framework. Datasets can be stored either on the \face hub or locally. HPC-FAIR, inheriting from \face, provides APIs for loading data from the hub using \face dataset identifiers or loading locally stored data by passing the dataset path.

% \textbf{POJ-104 dataset}~\cite{dataset:poj104} is derived from a pedagogical programming open judge (OJ) system that automatically evaluates the validity of submitted source code for specific problems by executing the code. This dataset is particularly useful for the code similarity task, as it consists of 104 problems, each accompanied by 500 student-written C/C++ programs that provide solutions to the respective problem.

% \textbf{DRB-ML-Dataset} is an OpenMP dataset, derived from \textbf{DataRaceBench}~\cite{dataset:dataracebench}, which contains 658 C /C++ OpenMP kernels.
% \textbf{DataRaceBench}~\cite{dataset:dataracebench} is an open-source benchmark suite developed to systematically and quantitatively evaluate the effectiveness of data race detection tools. DataRaceBench (DRB) focuses on parallel programs written in OpenMP and includes a set of micro-benchmark programs with or without data races to enable support for data race tool evaluation. The micro-benchmarks in DRB come from three main sources: manually written, extracted from real scientific applications, or automatically generated optimization variants.

%\todo[]{add dataset statistics table showing category and number of Qs in each category as well as examples \textcolor{blue}{this is done within the dataset part in the approach section}}
% \textbf{\crawldata} is an open-source benchmark composed of data from three resources: code crawled from GitHub, the OpenMP benchmark, and synthetic code. The design of the \crawldata dataset focuses on loops that can be parallelized with OpenMP. Each loop in the dataset is given a label indicating whether it is parallel and, if parallelizable, the corresponding OpenMP directive associated with the loop.
%\todo{Let's consistently use OpenMPQA as the name of the dataset, underscore in name is a bad choice for typing. \textcolor{blue}{We are using OMPQA right now}}
% \textbf{OMPQA} is created in this work for evaluating LLM models. We aim to create a comprehensive, diverse and balanced dataset and keep the following key principles to keep in mind while creating an OpenMP Q\&A dataset:
% \begin{enumerate}
%     \item Diversity of Topics: Include questions covering a wide range of topics related to OpenMP such as its basic syntax, data sharing, task scheduling, memory model, environment variables, thread affinity, synchronization, nested parallelism, and performance tuning. This will ensure that the model is trained on a diverse set of topics.
%     \item Multiple Difficulty Levels: The dataset should contain questions of varying difficulty levels, from beginner to advanced. This will make the model more adaptable in dealing with a wide range of queries, from basic to highly technical.
%     \item Different Types of Questions: Include a mix of "what is" (definition), "how to" (process), "why" (reasoning), and "which" (selection) questions. This variation will help the model recognize and correctly respond to different types of queries.
%     \item Real-world Problems: Include questions that are based on real-world problems or scenarios. These could be drawn from forums, textbooks, or academic papers. Real-world scenarios often involve complex, multi-faceted problems that challenge the model to synthesize information from different areas of knowledge.
%     \item Balanced Positive and Negative Examples: In addition to correct answers, also include incorrect ones (marked as such) to help the model learn what constitutes an incorrect response.
%     \item Data Privacy: The dataset should not contain any personally identifiable information or sensitive data. Ensure all data is anonymized and complies with relevant data protection laws and regulations.
% \end{enumerate}

%\todo{TODO: a table of question categories and sample Q\&A \textcolor{blue}{this is done within the dataset part in the approach section}}

% \begin{table}[h]
% \centering
% \begin{tabular}{|c|p{1in}|p{3in}|}
% \hline
%  Category &  Description & Examples    \\
% \hline
%  Basics& OpenMP concepts, quiz &  \\
% \hline
% Examples & example OpenMP code snippets & Give an example OpenMP code for computing PI  \\
% \hline
% Compilers & available OpenMP compilers, usage, internals &  \\
% \hline
%  Benchmarks & Available benchmarks, knowledge about individual benchmarks & What open-source benchmarks are available for measuring performance of OpenMP compilers? What benchmarks are available for evaluating data race detection tools?   \\
% \hline 
% \end{tabular}
% \caption{Categories of OpenMP Questions, with examples and Answers}
% % current collection of Q&A
% % https://docs.google.com/spreadsheets/d/1dqsd2v_lcRyLv8XiLIy00E_ezBiaivcSjZrMDqFeu_w/edit#gid=0
% \label{tab-omp-qa}
% \end{table}
%\todo{To lin: please add 10 Clang/LLVM compiler related questions}

%\todo{What is the common name for this task: Code similarity vs. Similarity Analysis vs. something else?  Please cite some previous paper and define the task when first mentioning this task name in Approach section, and in this Results section, we don't need to re-define the tasks. We just present results. \textit{similarity checking}}
\subsection{Code Similarity Analysis}
% \textbf{Task definition.}
% The definition of code similarity varies for different applications and datasets. As mentioned in section 3.2. The POJ-104 dataset defines similarity as code equivalent in functional. We process the POJ-104 dataset into tuples, in which each pair of code are assigned a label for similarity 
The code similarity task is designed to measure the syntactic and/or semantic similarity between two code snippets. Such analysis information can be beneficial in various scenarios such as plagiarism detection, code reuse and refactoring, bug detection and repair, licensing compliance, malware detection, and so on.
%The concept of code similarity can take on various forms depending on the specific application and dataset. 
% The code similarity task is designed to measure the semantic similarity between two code snippets. By evaluating the semantic relationship between the input code pair, the task aims to determine whether the codes share similar functionality or not. This task is formulated as a binary classification problem, with the output indicating whether the two code snippets are semantically similar (1) or not (0).

% \noindent
\textbf{Preparing Datasets and Ground Truth.}
%As mentioned in Section~\ref{sec-approach}, the POJ-104 dataset is a collection of code snippets, each associated with a label representing its specific functionality. To prepare the dataset for code similarity analysis, we first use \ourwork datasets API to load the POJ-104 dataset from the local drive and generate pairs of code snippets. 
%We use \ourwork datasets API to load 
Two datasets introduced in Section~\ref{sec-datasets}, POJ-104 and DRB-ML, are loaded through \ourwork's datasets API for this experiment.
For each pair of code snippets in the POJ-104 dataset, we assign a binary similarity label based on their functional labels.  A similarity label of 1 signifies that the snippet pair shares the same functional label and we assign a similarity score of 1. Otherwise, the label is 0. 
%Conversely, if the functional labels differ, we assign a similarity label of 0, denoting that the code snippets are functionally distinct or 'dissimilar'. 
We have processed the DRB-ML dataset using a similar methodology to generate code pairs and labels. The main difference is that the similarity ground truth for DRB-ML is derived from its own similarity score table~\cite{chen2022early}, providing a precise and reliable similarity measurement between code snippets in the dataset.
%It's noteworthy that the DRB-ML dataset is unbalanced after processing with its similarity loop-up table. With various metrics provided by \ourwork evaluator API, we strategically pick different metrics for evaluations with the two datasets. 

%\todo[]{add evaluation in listing 1}
\textbf{Inference Experiments and Evaluation.}
%For the rapid deployment of state-of-the-art models for the code similarity task, 
We employ \ourwork's code similarity pipeline to test various models. 
%This API enables users to effortlessly obtain inference results for input data using either the default model or a user-specified model that the pipeline supports.  
%\todo{\textcolor{blue}{added }Mention temperature settings as needed in the text}
The default model for this pipeline is CodeBERT. We additionally select four models from Table \ref{tab-modelExample} for evaluation: GraphCodeBERT, gpt-3.5-turbo, Dolly 2.0 (12B), and Cerebras-GPT (13B). 
We set the maximum token length for the model output to 256. This limits the verbosity of the model and keeps its responses concise. Additionally, we set the temperature parameter to 0 when applicable. For models like Dolly 2.0 that require positive temperature, we set the temperature to be $1\times 10^{-6}$. This setting ensures that the model's responses are consistent and deterministic, minimizing variability and uncertainty in its output. 

% In \ourwork, the handling of these models within the code similarity pipeline can be divided into two categories. On the one hand, CodeBert and GraphCodeBERT, which are language models specifically devised and trained on a variety of programming languages, are processed by comparing the code embeddings these models generate in order to assess code similarity. On the other hand, models like gpt-3.5-turbo, Dolly, and Cerebras, which are Large Language Models (LLMs) designed for a multitude of tasks and exhibit different parameter sizes and training data diversities, are managed differently. For these models, we directly input the codes as prompts and request answers. We use the default max\_token for each LLM and adjust the temperature to 0.5 for balanced output. This dual approach enables \ourwork to cater to a wide variety of models and tasks efficiently.

Within \ourwork, the approach of processing input code pairs depends on the type of the model employed. 
% Models like CodeBert and GraphCodeBert, which are specifically devised and trained on a variety of programming languages, directly embed the code pairs and pass to their model to generate a similarity prediction. 
Models like CodeBert and GraphCodeBert are specifically devised and trained on a variety of programming languages. We directly feed a pair of code snippets to generate a similarity prediction.
On the other hand, large language models like gpt-3.5-turbo, Dolly 2.0, and Cerebras-GPT are evaluated using the following prompt template: ``Code 1: \{...\} Code 2: \{...\}  Determine whether the two code snippets are similar. If the code snippets are similar, output 1; otherwise, output 0.''.


%\todo{\textcolor{yellow}{check embedding of LLMs; added under "Inference Experiments and Evaluation"} Not clear how LMs are used to generate similarity score: did you use questions to ask two code snippets or using their embeddings?}
%All the models are evaluated using their out-of-the-box configurations, with no specific tuning for the task. 

%\ourwork pipeline API loads models either from the \face model hub by passing the model identifier or from a locally saved checkpoint, as shown in listing \ref{ls:pipepline}. 
%The output from the pipelines can then be passed to the HPC-FAIR evaluator API. This API provides a quick and efficient means of model evaluation, including generating a leaderboard. % that allows users to compare and assess the performance of various models in a straightforward manner.
% To quickly deploy the cutting-edge models for the similarity task, users can use the \ourwork pipeline API by passing the model identifier or from locally saved checkpoint as illustrated in listing \ref{ls:pipepline}. The output of the pipelines can be passed to \ourwork evaluator API to generate leaderboard for a quick model evaluation. 

% \noindent
\textbf{Results.}
%We have evaluated two language  models and three large languae models for this task. 
%\todo{\textcolor{blue}{done }fix table ref}
The Code Similarity Analysis leaderboards generated using the two datasets are shown in Table~\ref{tab:hpctask1}.
%\todo{\textcolor{blue}{change right table to show F1 only, justifying by saving space}}
Notably, gpt-3.5-turbo demonstrates superior performance.
%in comprehending code snippets and comparing code similarity under various definitions. 
Two other models, StarChat-Alpha and Dolly 2.0, also exhibit commendable performance. 
Most large language models outperform traditional models (GraphCodeBERT and CodeBERT) that were specifically trained for code analysis. 
However, Cerebras-GPT struggled to comprehend the code and mostly returned arbitrary word tokens, indicating a lack of effective code understanding since it is mostly designed for natural language processing.
%\todo{Read Cerebras official tech report or paper: checking if code understanding tasks were used to evaluate it at all. PH please help this.}
%From these observations, it can be inferred that language models do possess the ability to understand programming languages and compare code similarity. In particular, powerful language models like gpt-3.5-turbo outperform traditional language models and lead the leaderboard for both datasets.

%In general, language models with a greater number of parameters perform better in this task. However, the diversity of the training data is also a crucial factor. Large language models that haven't been trained with code, like Cerebras, struggle with the code similarity checking task, emphasizing the importance of task-specific training data for effective model performance.





% HPC-FAIR evaluates the output of each model and generates the leaderboard of models for the similarity checking task. The evaluation results are shown in table \ref{tab:t1-drb}. \todo{What can we learn from the results. A table without explanation is not sufficient for a paper.}
%\todo{use bold numbers to highlight best scores for each column}
%\todo{The order of the rows: please consider order the best to the worst models for each tasks. This is how leaderboards usually show things. \textcolor{blue}{fixed}}


% \begin{table}[htbp]
%   \centering
%   \caption{Leaderboards for Code Similarity Analysis using the POJ-104 dataset (left) and the DRB dataset (right)}
% \label{tab:t1-drb}
%   \begin{minipage}{.5\linewidth}
%     \centering
% % \caption{Task 1 results with POJ-104 dataset}
% % \label{tab:t1-poj104}
% \begin{tabular}{|c|c|c|c|}
% \hline
% Model & Precision              & Recall                 & F1                     \\ \hline
% CodeBERT        & 51.5                   & 59.4                   & 55.2                   \\ \hline
% GraphCodeBERT   & 52.7                   & 60.3                   & 56.3                   \\ \hline
% ChatGPT 3.5     & 78.4                   & 74.2                   & 76.2                   \\ \hline
% Cerebras 13B    & 0                      & 0                      & 0                      \\ \hline
% Dolly 2.0       & 61.9                   & 61.3                   & 61.6
%     \\ \hline
% \end{tabular}
%   \end{minipage}%
%   \hspace{-1cm}
%   \begin{minipage}{.5\linewidth}
%     \centering
% % \caption{Task 1 results with DRB dataset}
% % \label{tab:t1-drb}
% \begin{tabular}{|c|c|}
% \hline
% Model & Accuracy                       \\ \hline
% CodeBERT        & 84.4                    \\ \hline
% GraphCodeBERT   & 85.1                             \\ \hline
% ChatGPT 3.5     & 88.7                          \\ \hline
% Cerebras 13B    & 0                               \\ \hline
% Dolly 2.0       & 76.9             
%     \\ \hline
% \end{tabular}
%   \end{minipage}
% \end{table}

% \setlength{\textfloatsep}{8pt}
\begin{table}[htbp]
\caption{Code Similarity Analysis Leaderboard: POJ-104 (left) and DRB-ML (right)}
\label{tab:hpctask1}
%  \centering
\begin{minipage}{.5\linewidth}
  % \caption{Leaderboard using POJ-104}
  % \label{tab:t1.1}
  \begin{tabular}{|c|c|c|c|}
  \hline
  \textbf{Model} & \textbf{Precision}              & \textbf{Recall}                 & \textbf{F1}                     \\ \hline
  \textbf{gpt-3.5-turbo}     & \textbf{78.4}                   & \textbf{74.2}                   & \textbf{76.2}         \\ \hline
  Dolly 2.0 12B  & 61.9                   & 61.3                   & 61.6                   \\ \hline
StarChat-Alpha  & 59.4                      & 56.2                      & 57.8            \\ \hline
GraphCodeBERT   & 52.7                   & 60.3                   & 56.3             \\ \hline
CodeBERT        & 51.5                   & 59.4                   & 55.2          \\ \hline
Cerebras-GPT 13B   & 0                      & 0                      & 0              \\ \hline        
\end{tabular}
\end{minipage}
\hspace{0.1cm}
\begin{minipage}{.5\linewidth}
  % \caption{Leaderboard using DRB-ML}
  % \label{tab:t1.2}
  \begin{tabular}{|c|c|c|c|}
\hline
\textbf{Model} & \textbf{Precision}              & \textbf{Recall}                 & \textbf{F1}                     \\ \hline
\textbf{gpt-3.5-turbo}     & \textbf{82.4}                   & \textbf{81.3}                   & \textbf{81.8}           \\ \hline
StarChat-Alpha  & 79.6                   & 77.4                   & 78.5                   \\ \hline
Dolly 2.0 12B  & 74.3                      & 73.2                      & 73.7            \\ \hline
GraphCodeBERT   & 79.4                   & 77.9                   & 78.6            \\ \hline
CodeBERT        & 76.9                   & 74.5                   & 75.7          \\ \hline
Cerebras-GPT 13B    & 0                      & 0                      & 0              \\ \hline            
\end{tabular}  
\end{minipage}
\end{table}

%\begin{table}[htbp]
%%  \centering
%\begin{minipage}{.5\linewidth}
%  \caption{Leaderboards for Code Similarity Analysis: POJ-104 dataset}
%  \label{tab:t1}
%%  \centering
%% \caption{Task 1 results with POJ-104 dataset}
%% \label{tab:t1-poj104}
%\begin{tabular}{|c|c|c|c|}
%\hline
%Model & Precision              & Recall                 & F1                     \\ \hline
%\textbf{gpt-3.5-turbo}     & \textbf{78.4}                   & \textbf{74.2}                   & \textbf{76.2}                   \\ %\hline
%Dolly 2.0 12B  & 61.9                   & 61.3                   & 61.6                   \\ \hline
%StarChat-Alpha   & 59.4                      & 56.2                      & 57.8            \\ \hline
%GraphCodeBERT   & 52.7                   & 60.3                   & 56.3             \\ \hline
%CodeBERT        & 51.5                   & 59.4                   & 55.2          \\ \hline
%Cerebras-GPT 13B    & 0                      & 0                      & 0              \\ \hline          
%\end{tabular}
%\end{minipage}%
%%  \hspace{-1cm}
%%   \begin{minipage}{.5\linewidth}
%%     \centering
%% % \caption{Task 1 results with DRB dataset}
%% % \label{tab:t1-drb}
%% \begin{tabular}{|c|c|}
%% \hline
%% Model & Accuracy                       \\ \hline
%% \textbf{gpt-3.5-turbo}     & \textbf{88.7}                    \\ \hline
%% GraphCodeBERT   & 85.1                             \\ \hline
%% CodeBERT        & 84.4                          \\ \hline
%% Dolly 2.0       & 76.9                           \\ \hline
%% Cerebras 13B    & 0           
%%     \\ \hline
%% \end{tabular}
%%   \end{minipage}
%%\end{table}
%% Model & Accuracy                       \\ \hline
%% CodeBERT        & 84.4                    \\ \hline
%% GraphCodeBERT   & 85.1                             \\ \hline
%% ChatGPT 3.5     & 88.7                          \\ \hline
%% Cerebras 13B    & 0                               \\ \hline
%% Dolly 2.0       & 76.9             
%
%%\todo{fix table}
%
%%\begin{table}[htbp]
%%  \centering
%\begin{minipage}{.5\linewidth}
%  \caption{Leaderboards for Code Similarity Analysis: DRB-ML dataset}
%\label{tab:t1}
%%    \centering
%% \caption{Task 1 results with POJ-104 dataset}
%% \label{tab:t1-poj104}
%\begin{tabular}{|c|c|c|c|}
%\hline
%Model & Precision              & Recall                 & F1                     \\ \hline
%\textbf{gpt-3.5-turbo}     & \textbf{82.4}                   & \textbf{81.3}                   & \textbf{81.8}                   \\ %\hline
%StarChat-Alpha  & 79.6                   & 77.4                   & 78.5                   \\ \hline
%Dolly 2.0   & 74.3                      & 73.2                      & 73.7            \\ \hline
%GraphCodeBERT   & 79.4                   & 77.9                   & 78.6            \\ \hline
%CodeBERT        & 76.9                   & 74.5                   & 75.7          \\ \hline
%Cerebras-GPT 13B    & 0                      & 0                      & 0              \\ \hline            
%\end{tabular}
%\end{minipage}%
%%  \hspace{-1cm}
%%   \begin{minipage}{.5\linewidth}
%%     \centering
%% % \caption{Task 1 results with DRB dataset}
%% % \label{tab:t1-drb}
%% \begin{tabular}{|c|c|}
%% \hline
%% Model & Accuracy                       \\ \hline
%% \textbf{gpt-3.5-turbo}     & \textbf{88.7}                    \\ \hline
%% GraphCodeBERT   & 85.1                             \\ \hline
%% CodeBERT        & 84.4                          \\ \hline
%% Dolly 2.0       & 76.9                           \\ \hline
%% Cerebras 13B    & 0           
%%     \\ \hline
%% \end{tabular}
%%   \end{minipage}
%\end{table}






% \begin{table}[]
% \centering
% \caption{Task 1 results with POJ-104 dataset}
% \label{tab:t1-poj104}
% \begin{tabular}{|c|c|c|c|}
% \hline
% Model-Tokenzier & Precision              & Recall                 & F1                     \\ \hline
% CodeBERT        & 51.5                   & 59.4                   & 55.2                   \\ \hline
% GraphCodeBERT   & 52.7                   & 60.3                   & 56.3                   \\ \hline
% ChatGPT 3.5     & 78.4                   & 74.2                   & 76.2                   \\ \hline
% Cerebras 13B    & 0                      & 0                      & 0                      \\ \hline
% Dolly 2.0       & 61.9                   & 61.3                   & 61.6
%     \\ \hline
% \end{tabular}
% \end{table}



% \begin{table}[]
% \centering
% \caption{Task 1 results with DRB dataset}
% \label{tab:t1-drb}
% \begin{tabular}{|c|c|}
% \hline
% Model-Tokenzier & Accuracy                       \\ \hline
% CodeBERT        & 84.4                    \\ \hline
% GraphCodeBERT   & 85.1                             \\ \hline
% ChatGPT 3.5     & 88.7                          \\ \hline
% Cerebras 13B    & 0                               \\ \hline
% Dolly 2.0       & 76.9             
%     \\ \hline
% \end{tabular}
% \end{table}
%\todo{Please add ranking columns to all leaderboards when possible}

%==================================================================
\subsection{Parallelism Detection}
\vspace{-5pt}
The parallelism detection task aims to identify parallelism opportunities within a given code snippet. We utilized two datasets, \crawldata  and DRB-ML introduced in Section~\ref{sec-datasets}, for the experiment. 
%Both datasets are created from previous research datasets to follow the format of \face such that it can be operated by \ourwork APIs. \todo{check if cited}

%\todo{\textcolor{blue}{added in listing 2 and explained in text} Some details are missing: need an concrete input and expected output. Does the input code snippet contain enough context information or just functions or loops? \textbf{Put input and output examples} \textit{added in the dataset part}}
%\todo{\textcolor{yellow}{token limits details not discussed, do we need to put it here? }\textcolor{blue}{explained in text with listing 2.}explain the input code snippets: 1. examples of previous work, understand not perfect; 2. input token limit; 3. }

% \noindent
\textbf{Preparing Datasets and Ground Truth.}
The \crawldata dataset is specifically designed for parallelism detection. 
%Each row in the dataset corresponds to a loop and labels indicating whether the loop is parallelizable and the type of OpenMP directive (if applicable). 
%This dataset includes widely-used benchmarks in the OpenMP community, such as NasBench ~\cite{} and Rodinia benchmark ~\cite{}. 
%After loading the \crawldata dataset, we extract only the label indicating whether the code is parallelizable, along with the data itself. Similar to the \crawldata dataset, a label of 1 represents parallelizable code, while a label of 0 indicates non-parallelizable code. 
Its existing labeling scheme allows us to prepare the data for binary classification models. 
Similarly, we prepared DRB-ML dataset with a label indicating whether each code snippet is parallelizable using OpenMP or not.  
%The parallelism detection task makes use of the entire DRB-ML dataset. However, as every piece of code in the DRB-ML dataset is parallelizable, it creates an imbalance for this task. Therefore, we utilize different metrics provided by HPC-FAIR to evaluate model performance.

It is worth noting that both datasets have undergone source code pre-processing steps, including comment removal and code snippet extraction. These steps are common practice~\cite{chen2023learning-dataset} to ensure that code snippets are small enough to be fed into language models with limited input token sequence sizes. However, the resulting code snippets may lose their context information, such as variable declarations. This is a serious limitation of language models with limited input sizes when applied to process large source files. 
% No space for the code
%Listing \ref{ls:parallel-data} shows an example code for the parallelism detection task from the DRB-ML dataset. 

% \begin{comment}
  

% \begin{listing}[h]
% \label{ls:parallel-data}
% \caption{An example code snippet from the DRB-ML dataset}
%     \begin{minted}[
% frame=lines,
% framesep=2mm,
% baselinestretch=1.2,
% fontsize=\scriptsize,
% linenos
% ]{c}
% for (c1 = 0; c1 <= ((nj + -1 < nm + -1 ? nj + -1 : nm + -1)); c1++) 
% {
%     for (c2 = 0; c2 <= ((nl + -1 < nm + -1 ? nl + -1 : nm + -1)); c2++) {
%         C[c1][c2] = ((double)c1) * (c2 + 3) / nl;
%         D[c1][c2] = ((double)c1) * (c2 + 2) / nk;
%     }
% ...// code omitted here
% }
% \end{minted}
% \end{listing}
% \end{comment}
%\todo{\textcolor{blue}{we have checked autoparbench. NPB and Rodinia is included in omp-serial dataset}a previous paper used NPB and Rodinia for testing auto parallelization. search: AutoParBench}

%\textbf{Inference Experiments and Evaluation} 
%HPC tasks like parallelism detection usually require specially designed models constructed by multiple machine learning layers and pre-trained with a large dataset. \ourwork aims to collect cutting-edge models and deploy them in pipelines according to specific tasks. 
%For the parallelism detection pipeline, two language models, augAST\cite{} and DeepSCC~\cite{} We reproduced previous works to train the models with the \crawldata training set and make the pretrained model accessable to users using \ourwork APIs. In this section, we evaluated three LLMs besides the two pre-trained language models. 


% \noindent
\textbf{Inference Experiments and Results.}
We selected six models to generate parallelism detection leaderboards. Four of them are introduced in Section~\ref{sec-background}.  They take the code snippets in a prompt template: ``As an OpenMP expert, you will analyze the given code snippet to determine if it can be parallelized. Code: \{...\}. Answer yes or no first:''. The other two are augAST~\cite{chen2023learning-dataset} and DeepSCC-based~\cite{harel2023learning}, which are pre-trained models using \crawldata's training dataset. We fed code snippets to these two models to directly obtain predicted labels. 

Table \ref{tab:t2} presents the resulting leaderboards. 
The highest F1 score reaches 93.9, indicating that LMs can be very effective for detecting parallelism. However, the datasets contain small-scale code snippets that are significantly simpler than real HPC codes. 
Again, gpt-3.5-turbo outperforms all other models overall, including specially trained models like augAST and DeepSCC.
AugAST performs better than gpt-3.5-turbo in terms of precision, suggesting it's more effective in predicting a positive class, which, in this case, is parallelizable code. 
%However, gpt-3.5-turbo scores better on recall and F1 measures, indicating it can effectively manage both true positive and false negative cases. 
%Language models trained with relatively smaller training datasets, like Dolly and Cerebras, do not perform as well as the fine-tuned models.
Finally, Cerebras-GPT did not perform well in this code analysis task. 

%gpt-3.5-turbo model has been trained with a massive amount of data. Although the specific details of the training data have not been disclosed, it's possible that it includes GitHub code. As a result, most of the data in the two datasets may have already been exposed to the gpt-3.5-turbo model. LLMs trained with a relatively smaller training data set, like Dolly and Cerebras, do not perform as well as the traditional pre-trained models.\todo{pre-trained vs. fine-tuned are different concept. I think you mean fine-tuned models for parallelism detection here.} Cerebras, in particular, struggles to comprehend the code in this task.
% & Model & Precision              & Recall                 & F1                     \\ \hline
% 1 & augAST        & 92.1                   & 82.4                   & 87.0                   \\ \hline
% 2& DeepSCC   & 82.7                   & 81.4                   & 82.0                   \\ \hline
% 3& ChatGPT 3.5     & 90.6                   & 89.3                   & 89.9                   \\ \hline
% 4& Cerebras 13B    & 0                      & 0                      & 0                      \\ \hline
% 5& Dolly 2.0       & 64.2                   & 63.7                   & 63.9
%\todo{Table 5: using F1 at least form the second table, if space permits, using the same metrics as the left table}

\begin{table}[htbp]
  \centering
  \caption{Parallelism Detection Leaderboards: \crawldata(left) and DRB-ML(right)}
\label{tab:t2}
  \begin{minipage}{.5\linewidth}
    \centering
    \resizebox{6cm}{!}{
\begin{tabular}{|c|c|c|c|c|}
\hline
\textbf{Model} & \textbf{Precision}              & \textbf{Recall}                 & \textbf{F1}                     \\ \hline
\textbf{gpt-3.5-turbo}     & 90.6                   & \textbf{89.3}                  & \textbf{89.9}                   \\ \hline
augAST        & \textbf{92.1}                   & 82.4                   & 87.0                   \\ \hline
DeepSCC   & 82.7                   & 81.4                   & 82.0                   \\ \hline
StarChat-Alpha   & 85.7                   & 68.2                  & 75.9                   \\ \hline
Dolly 2.0 12B      & 64.2                   & 63.7                   & 63.9                      \\ \hline
Cerebras-GPT 13B    & 0                      & 0                      & 0
    \\ \hline
\end{tabular} 
}
  \end{minipage}%
  \hspace{-1.13cm}
  \hspace{0.9cm}
   \begin{minipage}{.5\linewidth}
    \centering
    \resizebox{6cm}{!}
{
\begin{tabular}{|c|c|c|c|c|}
\hline
\textbf{Model} & \textbf{Precision}              & \textbf{Recall}                 & \textbf{F1}                     \\ \hline
\textbf{gpt-3.5-turbo}     & 90.0                   & \textbf{98.9}                  & \textbf{94.2}                   \\ \hline
augAST        & \textbf{91.4}                   & 72.3                   & 80.7                   \\ \hline
DeepSCC   & 80.4                   & 79.5                   & 79.9                   \\ \hline
StarChat-Alpha   & 81.9                   & 20.3                  & 32.5                   \\ \hline
Dolly 2.0 12B      & 40.0                   & 11.2                   & 2.17                      \\ \hline
Cerebras-GPT 13B    & 0                      & 0                      & 0
    \\ \hline
\end{tabular}
}
  \end{minipage}%
\end{table}



%==================================================================
\subsection{OpenMP Q\&A}
% current collection of Q&A
% https://docs.google.com/spreadsheets/d/1dqsd2v_lcRyLv8XiLIy00E_ezBiaivcSjZrMDqFeu_w/edit#gid=0
 In this experiment, we utilized \ourwork to evaluate the capabilities of several language models in answering questions related to OpenMP. This evaluation was conducted using the OMPQA dataset, introduced in Section \ref{sec-datasets}. 
 %each LM functions as an OpenMP Q\&A (Question and Answer) system. It is designed such that users can pose questions related to OpenMP, and in response, the LLM generates relevant answers. \ourwork allows for a practical examination of each LLM's understanding and application of OpenMP knowledge.

\textbf{Experiment Settings.} 
%As delineated in , OMPQA encompasses a diverse array of topics pertinent to OpenMP. Each entry in the OMPQA dataset has a corresponding answer that has been derived from the official OpenMP documentation, verified online resources, or expert human responses. 
%The prompt for this task, 

% We set the maximum token length for the model output to 256. This limits the verbosity of the model and keeps its responses concise. Additionally, we set the temperature parameter to 0. This setting ensures that the model's responses are consistent and deterministic, minimizing variability and uncertainty in its output. 
%\todo{\textcolor{blue}{done }We should set temperature to be 0 to have consistent answers. This is also the setting used by other evaluation papers. }

Each model receives the question in the following prompt template:
``You are an OpenMP expert. Please answer this question. Question: \{question\}''. Two metrics are selected to evaluate the quality of answers: the Bilingual Evaluation Understudy (BLEU) and ROUGE-L metrics. BLEU is a precision-oriented metric measuring the overlap of n-grams between the generated text and a set of reference texts. 
%It can consider multiple references but lacks a direct measure of semantic meaning. 
%NOT enough useful info.an extensively employed algorithm in the field of machine translation and natural language processing, developed to assess the quality of machine-translated text from one natural language to another. 
ROUGE-L (Recall-Oriented Understudy for Gisting Evaluation - Longest Common Subsequence) calculates the longest common subsequence (LCS) that appears in a left-to-right sequence in both the system-generated and reference summaries, thus providing a measure of the coherence and fluidity of the generated text.

\textbf{Results.} 
Table \ref{tab:t3-qa} displays the Q\&A leaderboard of several selected models. We additionally include the memory and execution time information. The experiments using gpt-3.5-turbo do not consume any local GPU memory since the model is invoked remotely through OpenAI's API.

Again, gpt-3.5-turbo unsurprisingly outperforms other LLMs, including StarChat-Alpha, Dolly 2.0, and Cerebras-GPT. However, the highest ROUGE-L F1 score of 0.259 indicates that all models have room for improvement in answering OpenMP questions. One reason is that many questions in OMPQA are open-ended and do not necessarily have a single correct answer. Also, the two metrics used do not sufficiently consider semantics.

% it becomes apparent that the augmentation of LLMs with documentation via LangChain enhances performance in the `Basic' and `Examples' categories but results in a downturn in the `Compiler' and `Benchmarks' categories. This discrepancy is likely due to the fact that our documentation does not encompass information pertinent to compiler and benchmark topics. Additionally, we assess the performance using CodeBLEU\cite{ren2020codebleu}, a novel metric specifically designed for the evaluation of code generation tasks. By our observation, the CodeBLEU scores for the 'Examples' category evaluate the generated code better, indicating that the HPC-specific metrics in LM4HPC help with the LLM evaluation for HPC tasks.

% Please add the following required packages to your document preamble:
% \usepackage{multirow}
% \usepackage[table,xcdraw]{xcolor}
% If you use beamer only pass "xcolor=table" option, i.e. \documentclass[xcolor=table]{beamer}
\begin{table}[]
\caption{Q\&A Leaderboard using the OMPQA dataset. The arrow indicates the performance changes when augmenting external knowledge base by LangChain.}
\label{tab:t3-qa}
\begin{tabular}{|c|c|c|c|c|ccc|}
\hline
 &
   &
   &
   &
   &
  \multicolumn{3}{c|}{\textbf{ROUGE}} \\ \cline{6-8} 
\multirow{-2}{*}{\textbf{Model}} &
  \multirow{-2}{*}{\begin{tabular}[c]{@{}c@{}}\textbf{\scriptsize CPU}\\ \scriptsize  Mem. (GB)\end{tabular}} &
  \multirow{-2}{*}{\begin{tabular}[c]{@{}c@{}}\textbf{\scriptsize GPU} \\ \scriptsize  Mem. (GB)\end{tabular}} &
  \multirow{-2}{*}{\textbf{Time}(s)} &
  \multirow{-2}{*}{\textbf{BLEU}} &
  \multicolumn{1}{c|}{Recall} &
  \multicolumn{1}{c|}{Prescision} &
  F1 \\ \hline
  \begin{tabular}[c]{@{}c@{}}gpt-3.5-turbo\\ + LangChain\end{tabular}
 &
  4.1 &
  0 &
  21.452 &
  \textbf{{0.147}$\uparrow$} &
  \multicolumn{1}{c|}{{0.347$\downarrow$}} &
  \multicolumn{1}{c|}{0.262$\downarrow$} &
  \textbf{0.259$\uparrow$} \\ \hline
gpt-3.5-turbo &
  4.2 &
  0 &
  12.749 &
  0.139 &
  \multicolumn{1}{c|}{\textbf{0.446}} &
  \multicolumn{1}{c|}{0.231} &
  0.257 \\ \hline
StarChat-Alpha &
  6.8 &
  18.9 &
  29.732 &
  0.082 &
  \multicolumn{1}{c|}{0.322} &
  \multicolumn{1}{c|}{0.149} &
  0.172 \\ \hline
  \begin{tabular}[c]{@{}c@{}}Dolly 2.0 12B\\ + LangChain\end{tabular}
  &
  27.4 &
  39.8 &
  7.217 &
  0.084$\uparrow$&
  \multicolumn{1}{c|}{0.228$\uparrow$} &
  \multicolumn{1}{c|}{0.232$\downarrow$} &
  0.182$\uparrow$ \\ \hline
Dolly 2.0 12B &
  27.1 &
  39.2 &
  8.147 &
  0.06 &
  \multicolumn{1}{c|}{0.208} &
  \multicolumn{1}{c|}{\textbf{0.312}} &
  0.148 \\ \hline
Cerebras-GPT 13B &
  52.6 &
  11.7 &
  590.763 &
  0.071 &
  \multicolumn{1}{c|}{0.319} &
  \multicolumn{1}{c|}{0.089} &
  0.112 \\ \hline
\end{tabular}
\end{table}


% Please add the following required packages to your document preamble:
% \usepackage{multirow}
% \usepackage[table,xcdraw]{xcolor}
% If you use beamer only pass "xcolor=table" option, i.e. \documentclass[xcolor=table]{beamer}
% \begin{table}[]
% \caption{Q\&A leaderboard by category}
% \label{tab:langchain-category}
% \begin{tabular}{|c|c|c|c|ccc|}
% \hline
%  &
%    &
%    &
%    &
%   \multicolumn{3}{c|}{\textbf{ROUGE}} \\ \cline{5-7} 
% \multirow{-2}{*}{\textbf{Category}} &
%   \multirow{-2}{*}{\textbf{Model}} &
%   \multirow{-2}{*}{\textbf{BLEU}} &
%   \multirow{-2}{*}{\textbf{CodeBLEU}} &
%   \multicolumn{1}{c|}{\textbf{Recall}} &
%   \multicolumn{1}{c|}{\textbf{Prescision}} &
%   \textbf{F1} \\ \hline
%  &
%   GPT3 &
%   0.175 &
%   NA &
%   \multicolumn{1}{c|}{0.476} &
%   \multicolumn{1}{c|}{0.251} &
%   0.329 \\ \cline{2-7} 
%  &
%   GPT3+LangChain &
%   \cellcolor[HTML]{32CB00}0.182 &
%   NA &
%   \multicolumn{1}{c|}{\cellcolor[HTML]{32CB00}0.498} &
%   \multicolumn{1}{c|}{\cellcolor[HTML]{32CB00}0.276} &
%   \cellcolor[HTML]{32CB00}0.355 \\ \cline{2-7} 
%  &
%   Dolly &
%   0.092 &
%   NA &
%   \multicolumn{1}{c|}{0.227} &
%   \multicolumn{1}{c|}{0.275} &
%   0.249 \\ \cline{2-7} 
% \multirow{-4}{*}{Basic} &
%   Dolly+LangChain &
%   \cellcolor[HTML]{32CB00}0.101 &
%   NA &
%   \multicolumn{1}{c|}{\cellcolor[HTML]{32CB00}0.371} &
%   \multicolumn{1}{c|}{\cellcolor[HTML]{32CB00}0.324} &
%   \cellcolor[HTML]{32CB00}0.346 \\ \hline
%  &
%   GPT3 &
%   0.152 &
%   0.165 &
%   \multicolumn{1}{c|}{0.371} &
%   \multicolumn{1}{c|}{0.247} &
%   0.297 \\ \cline{2-7} 
%  &
%   GPT3+LangChain &
%   \cellcolor[HTML]{32CB00}0.187 &
%   \cellcolor[HTML]{32CB00}0.266 &
%   \multicolumn{1}{c|}{\cellcolor[HTML]{FD6864}0.339} &
%   \multicolumn{1}{c|}{\cellcolor[HTML]{32CB00}0.337} &
%   \cellcolor[HTML]{32CB00}0.338 \\ \cline{2-7} 
%  &
%   Dolly &
%   0.103 &
%   0.131 &
%   \multicolumn{1}{c|}{0.245} &
%   \multicolumn{1}{c|}{0.194} &
%   0.217 \\ \cline{2-7} 
% \multirow{-4}{*}{Examples} &
%   Dolly+LangChain &
%   \cellcolor[HTML]{32CB00}0.112 &
%   \cellcolor[HTML]{32CB00}0.147 &
%   \multicolumn{1}{c|}{\cellcolor[HTML]{32CB00}0.287} &
%   \multicolumn{1}{c|}{\cellcolor[HTML]{FD6864}0.166} &
%   \cellcolor[HTML]{FD6864}0.21 \\ \hline
%  &
%   GPT3 &
%   0.147 &
%   NA &
%   \multicolumn{1}{c|}{0.375} &
%   \multicolumn{1}{c|}{0.367} &
%   0.371 \\ \cline{2-7} 
%  &
%   GPT3+LangChain &
%   \cellcolor[HTML]{FD6864}0.139 &
%   NA &
%   \multicolumn{1}{c|}{\cellcolor[HTML]{FD6864}0.319} &
%   \multicolumn{1}{c|}{\cellcolor[HTML]{FD6864}0.311} &
%   \cellcolor[HTML]{FD6864}0.315 \\ \cline{2-7} 
%  &
%   Dolly &
%   0.057 &
%   NA &
%   \multicolumn{1}{c|}{0.275} &
%   \multicolumn{1}{c|}{0.266} &
%   0.27 \\ \cline{2-7} 
% \multirow{-4}{*}{Compilers} &
%   Dolly+LangChain &
%   \cellcolor[HTML]{FD6864}0.046 &
%   NA &
%   \multicolumn{1}{c|}{\cellcolor[HTML]{FD6864}0.217} &
%   \multicolumn{1}{c|}{\cellcolor[HTML]{FD6864}0.246} &
%   \cellcolor[HTML]{FD6864}0.231 \\ \hline
%  &
%   GPT3 &
%   0.105 &
%   NA &
%   \multicolumn{1}{c|}{0.247} &
%   \multicolumn{1}{c|}{0.265} &
%   0.256 \\ \cline{2-7} 
%  &
%   GPT3+LangChain &
%   \cellcolor[HTML]{FD6864}0.094 &
%   NA &
%   \multicolumn{1}{c|}{\cellcolor[HTML]{FD6864}0.216} &
%   \multicolumn{1}{c|}{\cellcolor[HTML]{FD6864}0.254} &
%   \cellcolor[HTML]{FD6864}0.233 \\ \cline{2-7} 
%  &
%   Dolly &
%   0.086 &
%   NA &
%   \multicolumn{1}{c|}{0.208} &
%   \multicolumn{1}{c|}{0.312} &
%   0.25 \\ \cline{2-7} 
% \multirow{-4}{*}{Benchmarks} &
%   Dolly+LangChain &
%   \cellcolor[HTML]{FD6864}0.089 &
%   NA &
%   \multicolumn{1}{c|}{\cellcolor[HTML]{FD6864}0.204} &
%   \multicolumn{1}{c|}{\cellcolor[HTML]{FD6864}0.277} &
%   \cellcolor[HTML]{FD6864}0.235 \\ \hline
% \end{tabular}
% \end{table}

%One key feature of \ourwork is the capacity to integrate LangChain, which uses a vector store to allow question answering using additional user data. 
To enhance the capacity of large language models (LLMs) in accurately responding to OpenMP queries, we integrate the official OpenMP documentation into our process. 
%However, this integration can be challenged by the size limit of the documentation. To overcome this, 
We employ LangChain, a mechanism designed to efficiently store and retrieve language model embeddings, enabling us to accommodate large volumes of new data. To assess the efficacy of using LangChain to incorporate additional user data, we leverage its API to create a vector store. This vector store holds embeddings of text chunks derived from the OpenMP API Specification v5.2 (669 pages) and the OpenMP Application Programming Interface Examples v5.2.1 (575 pages). We then select two LangChain-supported models, GPT-3.5 and Dolly 2.0, to utilize the vector store as an additional resource for answering queries, thereby demonstrating the practical utility of our approach.
The results indicate slight improvements in both the BLEU and ROUGE-L F1 scores, increasing from 0.139 to 0.147 and from 0.257 to 0.259, respectively. However, there are mixed results for recall and precision metrics. gpt-3.5-turbo has a better recall, 0.446, compared to 0.347 of the Langchain approach.  %This suggests that the Vector Store API provided by \ourwork enables LLMs to leverage additional user databases, code, or documentation, enhancing their performance on HPC tasks.

Further, we examine the effectiveness of the LangChain approach across different question categories. When addressing `Basic' questions, the BLEU scores rise by 20.7\% and 9.8\% for gpt-3.5-turbo and Dolly 2.0, respectively. Additionally, we assess the LangChain performance using the CodeBLEU metric\cite{ren2020codebleu} for the `Examples' category, observing a score increase of 6.1\% and 12.2\% for gpt-3.5-turbo and Dolly 2.0, respectively. These observations indicate that augmenting LLMs with documentation via LangChain improves performance for `Basic' and `Examples' categories. However, for `Compilers' and `Benchmarks' categories, the performance of gpt-3.5-turbo and Dolly 2.0 diminishes when utilizing LangChain, recording an average BLEU score drop of 8.0\% and 7.9\%, respectively. This drop is likely because our documentation does not include information relevant to compiler and benchmark topics.



We also manually investigated the answers generated by the models. Overall, StarChat-Alpha delivers competitive results compared to GPT-3.5. It seems to be a good choice for people who want to use open-source language models based on our experiments. Research has indicated that GPT-4 surpasses GPT-3.5 in a variety of domains. However, as of now, API accessibility for GPT-4 has not been made publicly available. We plan to assess GPT-4's performance as soon as it becomes accessible and incorporate it into our framework if it benefits HPC tasks.
%\todo{\textcolor{yellow}{move to section3: Addressing input limit challenge. or supplying additional knowledge challenges or both. Longer input can include an additional context (knowledge).} Make sure there is text explaining what a model + LangChain means: best to have a diagram explaining how different pieces (model, prompt, vector store) work together. This can be part of the approach section, explaining the QA pipeline}
%\todo{\textcolor{blue}{mentioned we used only the text}limit of LangChain PDFLoader}
%\todo{\textcolor{blue}{added why gpt4 is not avaliable}add metrics designed for code output; add why GPT4 is no available}
%\todo{\textcolor{blue}{will upload with arXiv version}check some answers vs. reference answers: understanding why the scores are low in general}
%\todo{\textcolor{blue}{will upload with arXiv version} save and upload the generated answers/questions to GitHub}
%\todo[color=green]{create a new table for code category outputs}


% Please add the following required packages to your document preamble:
% \usepackage{multirow}


%\todo{\textcolor{blue}{For my experimetns, yes. It is normal since the deviation can be high for different tests for the same model}Dolly + LangChain is faster than using Dolly only?}




% If you use beamer only pass "xcolor=table" option, i.e. \documentclass[xcolor=table]{beamer}
% \begin{table}[h]
% \centering
% \caption{Q\&A Leaderboard using the OMPQA dataset}
% \label{tab:t3-qa}
% \begin{tabular}{|c|r|r|r|rrr|}
% \hline
% & \multicolumn{1}{c|}{} & \multicolumn{1}{c|}{} & \multicolumn{1}{c|}{} & \multicolumn{3}{c|}{ROUGE-L}                                       \\ \cline{5-7} 
% \multirow{-2}{*}{model} &
%   \multicolumn{1}{c|}{\multirow{-2}{*}{temp}} &
%   \multicolumn{1}{c|}{\multirow{-2}{*}{max\_token}} &
%   \multicolumn{1}{c|}{\multirow{-2}{*}{BLEU}} &
%   \multicolumn{1}{c|}{Recall} &
%   \multicolumn{1}{c|}{Precision} &
%   \multicolumn{1}{c|}{F1} \\ \hline
% \textbf{gpt-3.5-turbo}              & 0.5                   & 209                   & \textbf{0.086}               & \multicolumn{1}{r|}{0.4378} & \multicolumn{1}{r|}{0.1704} & \textbf{0.2190} \\ \hline
% Dolly 2.0 + LangChain             & 0.5                   & 209                   & 0.061                & \multicolumn{1}{r|}{0.2607} & \multicolumn{1}{r|}{0.2115} & 0.2004 \\ \hline

% Dolly 2.0 12B & 0.5  & 209        & 0.2737                & \multicolumn{1}{r|}{0.208} & \multicolumn{1}{r|}{0.312} & 0.148 \\ \hline
% Cerebras-GPT 13B           & 0.5                   & 209                   & 0.0703                & \multicolumn{1}{r|}{0.2661} & \multicolumn{1}{r|}{0.0816} & 0.1107 \\ \hline
% \end{tabular}
% \end{table}

% \todo{redo a few model +langChain experiments: should not use embeddings generated by model A to hook up with Model B. }


%\todo{add in table 6: space disk/memory; We should report space (memory, disk) complexity of the experiments, may not have enough space right now}

% Le
%\todo{\textcolor{blue}{will upload data and code after the submission }put all code to HPC-FAIR group on github's repo, share with collaborators}

\section{Related Work} %: Pei-Hung - 1
%Research and development in machine learning heavily rely on deep learning frameworks to provide building blocks and high-level programming interfaces in designing, training and validating a deep neural network. % Please avoid footnotes as much as possible
PyTorch and TensorFlow are the most popular frameworks, backed by Meta AI and Google, respectively.
Both frameworks are similar in many respects, including 1) providing low-level APIs for development, 2) supporting a rich collection of libraries, and 3) maintaining dedicated hubs - PyTorch Hub and TensorFlow Hub - for providing pre-trained ML models. 
% with dedicated hubs, PyTorch Hub and TensorFlow Hub, to provide pre-trained ML models.
%\todo{\textcolor{blue}{have replaced them with\face }Le: please make sure the paper use consistent Hugging Face or HuggingFace. }
\face is a large open-source community that builds tools to enable users to build, train, and deploy machine learning models based on open-source code and technologies.
\face is best known for its \texttt{Transformers} library, which exposes a collection of Python APIs to leverage state-of-the-art deep learning architectures for NLP tasks. 
With the goal to simplify end-to-end NLP tasks, \face \texttt{Transformers} offers a pipeline that performs all pre- and post-processing steps on the given input text data. The overall process of the model inference is encapsulated within these pipelines. With the pipeline, users only need to provide the input texts and the model for the task.  The remaining connections among a model and required pre- and post-processing steps are hidden within the pipeline implementation.

There were various research works and developments to improve the ML ecosystem to be Findable, Accessible, Interoperable, and Reproducible (FAIR).  These existing frameworks aim to make the models, datasets, or both FAIR.  Among these frameworks, HPCFAIR~\cite{verma2021hpcfair} focuses on providing support for model interoperability, search capabilities for datasets and models, and seamless integration into HPC workflows.  The work in \cite{yu2022seamless} extended this work to include support for interoperability across different framework implementations using ONNX and provision to retrain a model with transfer learning.
However, HPCFAIR framework relies on users to handle data pre- and post-processing. In comparison, LM4HPC is equipped to manage data processing within the pipeline design and generate leaderboards for supported HPC tasks. 
%\todo{To Murali: how to differentiate this paper from previous ones? }

%LangChain~\cite{langchain2023} is a framework built to support language models with a goal to provide data-aware connection, between a language model and sources of data, and allow a language model to interact with its environment.
%The core of LangChain is to chain different components to create a effective usage of the LLMs.  LM4HPC APIs proposed in this work can leverage the features from LangChain to enable more sophisticated tasks with LLMs.

General LLMs are trained with data covering general knowledge and information that is usually collected from public domains. 
Domain-specific datasets can be collected for the training of a specialized model or for the fine-tuning of a general-purpose model.
MedQA\cite{jin2020disease} is an example of domain-specific datasets collecting question-answer pairs and textbooks  from professional medical board exams.
ExeBench~\cite{exebench}, another domain-specific dataset for tasks in compilation and software engineering, contains millions of runnable and representative C functions collected from GitHub.
In addition to collecting existing data, ML research has started to automate dataset creation with assistance from the LLMs.  
The developers of LaMini-LM~\cite{wu2023lamini} develop a large set of 2.58M instruction and response pairs based on both existing and newly-generated instructions.  A handful of seed examples from the existing LLM prompts and 2.2M categories from Wikipedia from existing  are submitted to the \texttt{gpt-3.5-turbo} to generate relevant instructions.  Similarly, the responses for the generated instructions are also generated by the \texttt{gpt-3.5-turbo}.

%Mention a few generic and domain-specific QA datasets (e.g. MedQA dataset) and creation methods: 
%https://arxiv.org/abs/2304.14402 TODO

%\todo{Mention a few generic and domain-specific QA datasets (e.g. MedQA dataset) and creation methods: https://arxiv.org/abs/2304.14402}
\label{sec-related-work}
% Pei-Hung

\section{Conclusion} %: Leo - 0.25
In this paper, we presented our efforts to facilitate the application of language models for tasks specific to High-Performance Computing. We have developed the LM4HPC framework to encompass and expose relevant machine learning components via corresponding APIs.
%We discuss several challenges we encountered, along with the solutions we implemented. Using three representative tasks, we evaluated our prototype framework and APIs in the context of generating leaderboards for select language models. 
Our experimental findings suggest that GPT-3 performs competitively, despite not being specifically designed for HPC tasks. However, there is significant room for improvement in answering OpenMP questions. Furthermore, the input size limitation of language models adds complexity to certain tasks, such as parallelism detection. Finally, an obstacle to advancing the application of language models for HPC tasks is the absence of HPC-specific training and evaluation datasets.

Looking ahead, our future work will explore automated approaches to generating HPC-specific datasets. We intend to enhance LM4HPC's capabilities to support the fine-tuning of models for HPC-related tasks, including those related to the Message Passing Interface (MPI), and to provide performance analysis and optimization suggestions.

% \todo{{1. StarCoder+LangChain; 2. API table; 3. OMP Documentation version; 4. Temperature 0; 5. Appendix}}
\label{sec-conclusion}

\section*{Acknowledgement}
Prepared by LLNL under Contract DE-AC52-07NA27344 (LLNL-CONF-849438) and supported by the U.S. Department of Energy, Office of Science, Advanced Scientific Computing Research.

%\todo{\textcolor{blue}{go with arXiv version} add An Appendix showing how to reproduce all results of this paper: scripts, models, datasets, example inputs and outputs of essential steps.}
%\listoftodos


% \todo{To save space, remove unessential citations. 20-25 citations are sufficient. }
\bibliographystyle{splncs04}
\bibliography{main.bib}% common bib file

\end{document}




